% ============================================================
% MODELO ESTRUTURAL
% ============================================================

% -------------------------------------------------------
\section{Modelo Estrutural}
% -------------------------------------------------------

\begin{frame}[plain]
    \vfill
    \centering
    \begin{beamercolorbox}[sep=12pt,center,shadow=false,rounded=true]{title}
        {\Large\bfseries Modelo Estrutural}\\[6pt]
        {\normalsize Equilíbrio espacial com comércio e migração custosos}
    \end{beamercolorbox}
    \vfill
\end{frame}

% -------------------------------------------------------
\begin{frame}{Estrutura do modelo}
    Modelo de equilíbrio geral espacial com $N$ localidades, combinando
    \textbf{comércio custoso} (Eaton e Kortum, 2002) e \textbf{migração custosa} (Monte et al., 2018).

    \vspace{0.3cm}

    \begin{columns}[T]
        \begin{column}{0.48\textwidth}
            \textbf{Blocos do lado da demanda}
            \begin{itemize}
                \item Utilidade indireta do consumidor.
                \item Decisão de migração com custos $\kappa_{odt}$.
            \end{itemize}

            \vspace{0.25cm}

            \textbf{Blocos do lado da oferta}
            \begin{itemize}
                \item Comércio com produtividade Fréchet e custos iceberg $\tau_{odt}$.
                \item Determinação de salários (mercado de bens).
                \item Determinação de aluguéis (mercado de habitação).
            \end{itemize}
        \end{column}
        \begin{column}{0.48\textwidth}
            \textbf{Objetivo final}
            \begin{itemize}
                \item Efeitos de bem-estar via mudanças em $\kappa$ e $\tau$ induzidas pelas estradas.
            \end{itemize}

            \vspace{0.25cm}

            \textbf{Custos de fricção}
            \begin{itemize}
                \item $\kappa_{odt} \geq 1$: custo de migrar de $o$ para $d$.
                \item $\tau_{odt} \geq 1$: custo iceberg de comerciar de $o$ para $d$.
                \item Ambos dependem do \textbf{tempo de viagem rodoviário}.
            \end{itemize}
        \end{column}
    \end{columns}
\end{frame}

% -------------------------------------------------------
\begin{frame}{Utilidade do consumidor}
    Condicional a viver em $d$, o indivíduo consome um agregado CES de bens e habitação
    segundo uma utilidade Cobb-Douglas. A \textbf{utilidade indireta} é:

    \begin{equation}\label{eq:utility}
        V_{dt} \;\propto\; B_{dt}\; P_{dt}^{-(1-\alpha)}\; r_{dt}^{-\alpha}\; w_{dt}
    \end{equation}

    \vspace{0.1cm}

    \begin{columns}[T]
        \begin{column}{0.48\textwidth}
            \textbf{Variáveis}
            \begin{itemize}\small
                \item $B_{dt}$: amenidade de $d$ no tempo $t$.
                \item $P_{dt}$: índice de preços de bens tradables.
                \item $r_{dt}$: aluguel (custo de habitação).
                \item $w_{dt}$: salário em $d$.
                \item $\alpha$: participação da habitação na renda.
            \end{itemize}
        \end{column}
        \begin{column}{0.48\textwidth}
            \textbf{Intuição}
            \begin{itemize}\small
                \item $V_{dt}$ é maior onde salários são altos, preços baixos e aluguéis baixos.
                \item Estradas $\downarrow$ $P_{dt}$ (bens mais baratos) \textbf{e} $\downarrow$ $\kappa_{odt}$ (migrar mais fácil) $\Rightarrow$ dois canais de ganho.
            \end{itemize}
        \end{column}
    \end{columns}
\end{frame}

% -------------------------------------------------------
\begin{frame}{Decisão de migração}
    Cada indivíduo $i$ em $o$ escolhe $d$ para maximizar $V_{dt}\,\kappa_{odt}^{-1}\,\xi_{dt}(i)$,
    onde $\xi_{dt}(i)$ é um choque idiossincrático. A fração de pessoas de $o$ que migra para $d$ é:

    \begin{equation}\label{eq:migration}
        \pi_{odt}^{\text{mig}} = \frac{V_{dt}^{\,\epsilon}\;\kappa_{odt}^{-\epsilon}}{\Phi_{ot}},
        \qquad
        \Phi_{ot} = \sum_{d'} V_{d't}^{\,\epsilon}\;\kappa_{od't}^{-\epsilon}
    \end{equation}

    \vspace{0.4cm}

    \begin{columns}[T]
        \begin{column}{0.48\textwidth}
            \textbf{Parâmetros e variáveis}
            \begin{itemize}\small
                \item $\epsilon$: elasticidade de migração a salários reais.
                \item $\kappa_{odt}$: custo de migrar de $o$ para $d$.
                \item $\Phi_{ot}$: atratividade agregada dos destinos para quem parte de $o$.
            \end{itemize}
        \end{column}
        \begin{column}{0.48\textwidth}
            \textbf{Intuição}
            \begin{itemize}\small
                \item Gravidade: migração $\uparrow$ com $V_{dt}$ e $\downarrow$ com $\kappa_{odt}$.
                \item Estradas $\downarrow$ $\kappa_{odt}$ $\Rightarrow$ $\pi_{odt}^{\text{mig}}$ $\uparrow$.
            \end{itemize}
        \end{column}
    \end{columns}
\end{frame}

% -------------------------------------------------------
\begin{frame}{Comércio e determinação de preços}
    Produtividade de $o$ no bem $j$: $\zeta_{ot}(j) \sim$ Fréchet($\theta$, $A_{ot}$). Com custo iceberg $\tau_{odt}$,
    $d$ importa de quem oferece o menor preço. A \textbf{fração de bens} importada de $o$ e o \textbf{índice de preços} são:

    \begin{equation}\label{eq:trade}
        \pi_{odt}^{\text{trade}} = \frac{A_{ot}\,w_{ot}^{-\theta}\,\tau_{odt}^{-\theta}}{\Theta_{dt}},
        \quad
        \Theta_{dt} = \sum_{o'} A_{o't}\,w_{o't}^{-\theta}\,\tau_{o'dt}^{-\theta},
        \quad
        P_{dt} \propto \Theta_{dt}^{-1/\theta}
    \end{equation}

    \vspace{0.1cm}

    \begin{columns}[T]
        \begin{column}{0.52\textwidth}
            \textbf{Variáveis}
            \begin{itemize}\small
                \item $A_{ot}$: produtividade média de $o$.
                \item $\theta$: elasticidade de comércio (vantagem comparativa).
                \item $\tau_{odt}$: custo iceberg ($\geq 1$) de $o$ para $d$.
                \item $\Theta_{dt}$: acesso de $d$ a fontes baratas/produtivas.
            \end{itemize}
        \end{column}
        \begin{column}{0.44\textwidth}
            \textbf{Intuição}
            \begin{itemize}\small
                \item Estrutura análoga a \eqref{eq:migration}: gravidade com fricções.
                \item Estradas $\downarrow$ $\tau_{odt}$ $\Rightarrow$ $\Theta_{dt}$ $\uparrow$ $\Rightarrow$ $P_{dt}$ $\downarrow$.
                \item Locais bem conectados têm preços menores.
            \end{itemize}
        \end{column}
    \end{columns}
\end{frame}

% -------------------------------------------------------
\begin{frame}{Equilíbrio de salários e aluguéis}
    Salários -- equilíbrio no mercado de bens:

    \begin{equation}\label{eq:wages}
        w_{ot}\,L_{ot} \;=\; \sum_{d}\;\pi_{odt}^{\text{trade}}\;w_{dt}\,L_{dt}
        \qquad \forall\; o \in \{1,\ldots,N\}
    \end{equation}

    \small Renda de $o$ = soma das exportações para todos os destinos. Sistema de $N$ equações que determina salários implicitamente.

    \vspace{0.2cm}

    Aluguéis -- equilíbrio no mercado de habitação:

    \begin{equation}\label{eq:rent}
        r_{dt} \;=\; (\nu_t)^{\frac{1}{1+\epsilon_r}} \left(\alpha\, w_{dt}\,L_{dt}\right)^{\!\frac{\epsilon_r}{1+\epsilon_r}}
    \end{equation}

    \small
    \begin{itemize}
        \item $\nu_t$: custo de construção; $\epsilon_r$: elasticidade inversa da oferta de habitação.
        \item Aluguel $\uparrow$ com $L_{dt}$ e $w_{dt}$ -- aglomeração pressiona moradia.
        \item Oferta inelástica: aluguéis absorvem ganhos salariais, comprimindo bem-estar líquido.
    \end{itemize}
\end{frame}

% -------------------------------------------------------
\begin{frame}{Efeitos de bem-estar}
    O bem-estar médio de origem $o$ é proporcional a $\Phi_{ot}^{1/\epsilon}$.
    A \textbf{variação agregada de bem-estar} é uma média ponderada das variações de oportunidade:

    \begin{align}
        \widehat{\text{bem-estar}}_t &= \sum_o \pi_{o,t-1}^{\text{mig}}\;\widehat{\Phi}_{ot}^{\,1/\epsilon} \label{eq:welfare} \\[3pt]
        \widehat{\Phi}_{ot} &= \sum_d \pi_{odt-1}^{\text{mig}}\;\widehat{V}_{dt}^{\,\epsilon}\;\widehat{\kappa}_{od}^{-\epsilon} \label{eq:phi} \\[3pt]
        \widehat{V}_{dt}   &= \widehat{P}_{dt}^{-(1-\alpha)}\;\widehat{r}_{dt}^{-\alpha}\;\widehat{w}_{dt} \label{eq:vhat}
    \end{align}

    \begin{itemize}\small
        \item \textbf{``Hat''}: $\widehat{X}_t = X_t / X_{t-1}$ -- razão entre situação nova e anterior.
        \item $\widehat{\Phi}_{ot}$ muda quando destinos ficam mais atrativos ($\widehat{V}_{dt}$ $\uparrow$) ou migrar fica mais fácil ($\widehat{\kappa}_{od}$ $\downarrow$).
        \item $\widehat{V}_{dt}$ melhora se salários $\uparrow$ ou preços/aluguéis $\downarrow$ -- exatamente os canais das estradas.
        \item O modelo separa: quanto do ganho vem de $\widehat{V}$ (comércio) \textit{vs.}\ $\widehat{\kappa}$ (migração).
    \end{itemize}
\end{frame}

% -------------------------------------------------------
\begin{frame}{Parâmetros estruturais estimados}
    \begin{columns}[T]
        \begin{column}{0.52\textwidth}
            \footnotesize\renewcommand{\arraystretch}{1.35}
            \begin{tabular}{llp{3cm}}
                \toprule
                Parâmetro & Valor & Fonte \\
                \midrule
                $\theta$ (comércio) & 4,0 & Simonovska-Waugh (2014) \\
                $\epsilon$ (migração) & 4,5 & Estimado -- Seção 4 \\
                $1/\epsilon_r$ (habitação) & 1,7 & Estimado -- Seção 4 \\
                $\alpha$ (hab.) & calibrado & Dados de consumo \\
                \bottomrule
            \end{tabular}
        \end{column}
        \begin{column}{0.44\textwidth}
            \vspace{-0.1cm}
            \normalsize\textbf{Intuição dos parâmetros}
            \begin{itemize}\small
                \item \textbf{$\theta = 4$}: comércio responde fortemente a reduções de custo.
                \item \textbf{$\epsilon = 4{,}5$}: heterogeneidade de gostos moderada -- migração responde, mas com fricção.
                \item \textbf{$1/\epsilon_r = 1{,}7$}: oferta inelástica -- aluguéis comprimem bem-estar.
            \end{itemize}
        \end{column}
    \end{columns}

    {\tiny\textit{$\epsilon$ e $1/\epsilon_r$: estimados via instrumento Bartik (Seção 4).}}
\end{frame}
